\documentclass[12pt, openright,oneside,a4paper,english,portuguese]{abntex2}

\include{Arquivos/pacotes}
\include{Dados}
\selectlanguage{portuguese}
\begin{document}
% CAPA
\include{textos/preambulo/capa}
%\begin{fichacatalografica}
    % O Arquivo abaixo precisa ser substituído pela ficha catalográfica emitida pela Biblioteca.
 %   \includepdf{Arquivos/Ficha_Catalografica.pdf}
%\end{fichacatalografica}
% -------------------------------------------------
% ELEMENTOS PRÉ-Textuais
% Os nomes são auto-explicativos
% -------------------------------------------------
\include{textos/preambulo/declaracaoAutoria}
\include{textos/dedicatoria}
\include{textos/agradecimentos}
%\include{textos/epigrafe}
\include{textos/resumo}
\include{textos/abstract}
\include{textos/preambulo/ListaFigurasTabelas}
\include{textos/listaSiglas}
\include{textos/listaSimbolos}
% ------------------------------------
% SUMARIO
% ------------------------------------
\pdfbookmark[0]{\contentsname}{toc}
\tableofcontents*
\cleardoublepage
\pagenumbering{arabic}
%-------------------------------------------------
% CAPÍTULOS
%-------------------------------------------------
%Exemplos no Latex, comente a linha abaixo na hora de compilar a versão final
%\include{capitulos/exemplo}

%Capítulo 1- Introdução
\include{capitulos/Capitulo1}
%Capítulo 2- Revisão da Literatura
\include{capitulos/Capitulo2}
%Capítulo 3- Propostas do Trabalho
\include{capitulos/Capitulo3}
%Capítulo 5 - Conclusão
\include{capitulos/conclusao}
%Capítulo 4- Recomendações
\include{capitulos/Recomendacoes}
% -------------------------------------------------------
% ELEMENTOS PÓS-TEXTUAIS
% -------------------------------------------------------
\bibliography{Bibliografia.bib}
\include{textos/anexos}
\include{textos/apendices}
\end{document}
